% !TeX root = ../高等代数习题（谢启鸿）笔记.tex
% 由主文档通过 \input 载入；请编译主文档。
\setcounter{chapter}{3} % 保留原书章号。
\chapter{线性映射}
\label{chap:xie-ex-linear-maps}

% 依据《高等代数习题》（谢启鸿）第 4 章整理：高代习题-4.pdf，
% PDF 第 1--42 页（原书第 192--233 页），覆盖 §§4.1--4.8。
% 将重复例题的不同解法合并；原书题号仅在注释或正文中说明，不作为本笔记编号。
线性映射把向量空间的运算与矩阵的运算联系起来。本章围绕两条相互配合的思路展开：一条是在空间中选基、取补空间、研究像与核，另一条是把映射写成矩阵，利用秩、相似与分块运算。做题时先明确定义域、陪域和数域，再判断适合在哪一种语言中推导。以下 $K$ 表示数域，即 $\mathbb C$ 的子域；没有另作说明时，所讨论的空间均为 $K$ 上的有限维线性空间，$\mathcal L(V,U)$ 表示 $V\to U$ 的全体线性映射，$\mathcal L(V)=\mathcal L(V,V)$。无限维情形将单独指出。$K^n$ 统一采用列向量，$M_{m\times n}(K)$ 表示 $m\times n$ 矩阵空间，$\operatorname{GL}_n(K)$ 表示可逆 $n$ 阶矩阵全体。

\section{基本概念}
\label{sec:xie-ex-linear-basics}

\subsection{线性映射及其运算}

\begin{definition}
\label{def:xie-ex-linear-map}
映射 $T:V\to U$ 若满足
\[
    T(ax+by)=aT(x)+bT(y)
    \qquad(x,y\in V,\ a,b\in K),
\]
称为\emph{线性映射}（linear map）。$V$ 到自身的线性映射称为\emph{线性变换}（linear transformation）；双射的线性映射称为\emph{线性同构}（linear isomorphism），$V$ 到自身的线性同构称为\emph{自同构}（automorphism）。
\end{definition}

线性性等价于保持加法与数乘，因而线性映射保持零向量和任意有限线性组合。特别地，$T(0)=0$ 是必要条件，但不是充分条件。对 $S,T\in\mathcal L(V,U)$、$c\in K$，规定 $(S+T)(x)=S(x)+T(x)$、$(cT)(x)=cT(x)$，便使 $\mathcal L(V,U)$ 成为线性空间。若 $T:V\to U$、$R:U\to W$，规定 $RT=R\circ T$，即 $(RT)(x)=R(T(x))$；复合保持线性，满足结合律、分配律和与数乘的相容性。于是 $\mathcal L(V)$ 以复合为乘法成为含幺结合代数，单位元是恒等变换 $I_V$，但乘法一般不交换。\emph[caution]{复合时先作用右侧映射，只有定义域与陪域匹配时才能相乘。}

\subsection{表示矩阵与换基}

取 $V$ 的有序基 $\mathcal E=(e_1,\ldots,e_n)$ 和 $U$ 的有序基 $\mathcal F=(f_1,\ldots,f_m)$。若
\[
    T(e_j)=\sum_{i=1}^{m}a_{ij}f_i\qquad(1\leq j\leq n),
\]
则 $A=(a_{ij})\in M_{m\times n}(K)$ 称为 $T$ 的\emph{表示矩阵}（representing matrix），记为 $[T]_{\mathcal F\leftarrow\mathcal E}$；第 $j$ 列是 $T(e_j)$ 在陪域基下的坐标。若 $[x]_{\mathcal E}$ 表示坐标列向量，则
\begin{equation}
\label{eq:xie-ex-coordinate-map}
    [T(x)]_{\mathcal F}=A[x]_{\mathcal E}.
\end{equation}
因此取定两组基后，映射 $T\mapsto A$ 是 $\mathcal L(V,U)\to M_{m\times n}(K)$ 的线性同构；复合映射的矩阵是相应矩阵的乘积。对于 $V$ 上的变换，讨论相似时定义域与陪域使用同一组基。

若新基满足 $\mathcal E'=\mathcal E P$、$\mathcal F'=\mathcal F Q$，其中 $P,Q$ 可逆，则坐标关系为 $[x]_{\mathcal E}=P[x]_{\mathcal E'}$、$[y]_{\mathcal F}=Q[y]_{\mathcal F'}$。代入式~\eqref{eq:xie-ex-coordinate-map} 得
\begin{equation}
\label{eq:xie-ex-general-change-basis}
    [T]_{\mathcal F'\leftarrow\mathcal E'}=Q^{-1}AP.
\end{equation}
特别地，同一变换在两组基下的矩阵满足 $B=P^{-1}AP$，称 $A,B$ \emph{相似}（similar）。一般映射的两侧换基相互独立，对应矩阵相抵；同一变换的同基换基对应相似，不能混用这两种关系。

\subsection{像、核与维数公式}

\begin{definition}
\label{def:xie-ex-image-kernel}
线性映射 $T:V\to U$ 的\emph{像空间}（image）和\emph{核空间}（kernel）分别是
\[
    \operatorname{Im}T=\{T(x):x\in V\}\subseteq U,
    \qquad \ker T=\{x\in V:T(x)=0\}\subseteq V.
\]
它们都是子空间。$\operatorname{rank}T=\dim\operatorname{Im}T$ 称为秩，$\dim\ker T$ 称为零度（nullity）。
\end{definition}

\begin{theorem}
\label{thm:xie-ex-rank-nullity}
若 $\dim V=n$，则
\begin{equation}
\label{eq:xie-ex-rank-nullity}
    \dim\ker T+\dim\operatorname{Im}T=n.
\end{equation}
若 $A$ 是 $T$ 的表示矩阵，则 $\operatorname{rank}T=\operatorname{rank}A$。因此 $T$ 单射当且仅当 $\operatorname{rank}A=n$，$T$ 满射当且仅当 $\operatorname{rank}A=\dim U$；当 $\dim U=n$ 时，单射、满射、同构和 $A$ 可逆相互等价。
\end{theorem}

\begin{proof}
取 $\ker T$ 的基 $k_1,\ldots,k_s$，扩充为 $V$ 的基 $k_1,\ldots,k_s,v_1,\ldots,v_{n-s}$。任意向量的像均为 $T(v_1),\ldots,T(v_{n-s})$ 的线性组合。若 $\sum_j c_jT(v_j)=0$，则 $\sum_jc_jv_j\in\ker T$，由上述基的线性无关性知所有 $c_j=0$，故这些像组成 $\operatorname{Im}T$ 的基，得到维数公式。坐标同构把 $\ker T$ 对应到 $Ax=0$ 的解空间，把 $\operatorname{Im}T$ 对应到 $A$ 的列空间，故两种秩相等。最后，线性映射单射当且仅当核为零，满射当且仅当像等于陪域，其余结论由维数公式得到。
\end{proof}

\subsection{不变子空间与分块形式}

\begin{definition}
\label{def:xie-ex-invariant-subspace}
对 $T\in\mathcal L(V)$，若子空间 $W\subseteq V$ 满足 $T(W)\subseteq W$，称 $W$ 为 $T$ 的\emph{不变子空间}（invariant subspace）。$\{0\}$ 和 $V$ 称为平凡不变子空间。
\end{definition}

若 $e_1,\ldots,e_r$ 是 $W$ 的基，并扩充为 $V$ 的基，则 $W$ 不变当且仅当 $T$ 的矩阵具有形式
\begin{equation}
\label{eq:xie-ex-invariant-block}
    [T]=\begin{pmatrix}A_{11}&A_{12}\\0&A_{22}\end{pmatrix},
\end{equation}
其中 $A_{11}$ 是限制变换 $T|_W$ 的矩阵。理由是前 $r$ 个基向量的像在 $W$ 中，其后 $n-r$ 个坐标均为零。若 $V=W_1\oplus\cdots\oplus W_s$ 且每个 $W_i$ 都不变，拼接各子空间的基便得到分块对角矩阵。反过来，分块对角矩阵中各坐标块张成的子空间都是不变子空间。

\section{线性映射及其运算}
\label{sec:xie-ex-linear-construction}

\subsection{在线性无关组与直和上构造映射}

% 原书例 4.1、4.2、4.4。
\begin{theorem}
\label{thm:xie-ex-linear-extension}
给定 $V$ 的一组基 $e_1,\ldots,e_n$ 和 $U$ 中任意向量 $u_1,\ldots,u_n$，存在唯一线性映射 $T:V\to U$ 使 $T(e_i)=u_i$。若 $V=V_1\oplus\cdots\oplus V_s$，给定 $T_i\in\mathcal L(V_i,U)$，也存在唯一线性映射 $T:V\to U$，使 $T|_{V_i}=T_i$。
\end{theorem}

\begin{proof}
对 $x=\sum_i a_ie_i$ 定义 $T(x)=\sum_i a_iu_i$。坐标的唯一性保证定义良好，线性组合的运算保证线性；任何满足基向量取值要求的线性映射都必须遵守此公式，故唯一。对直和情形，将唯一分解 $x=x_1+\cdots+x_s$ 映为 $\sum_iT_i(x_i)$，同样验证即可。
\end{proof}

此结论称为\emph{线性扩张定理}（linear extension theorem）。给定子空间 $W\subseteq V$ 上的线性映射时，选补空间 $V=W\oplus W'$，在 $W'$ 上任意指定线性映射即可扩张到 $V$；常取零映射以简化构造。两个映射只要在一组基上的值相同，便在整个空间上相等。若给定的是线性相关的生成组，则其像不能任意指定，而必须保持原有的全部线性关系，详见第~\ref{sec:xie-ex-linear-training} 节。

\subsection{单侧逆、广义逆与因子分解}

% 原书例 4.3、4.5--4.7；合并几何与矩阵两种证明。
\begin{proposition}
\label{prop:xie-ex-splitting-inverses}
设 $T:V\to U$ 线性。$T$ 单射当且仅当存在 $L:U\to V$ 使 $LT=I_V$；$T$ 满射当且仅当存在 $R:U\to V$ 使 $TR=I_U$。此外，总存在 $G:U\to V$ 使 $TGT=T$。
\end{proposition}

\begin{proof}
$LT=I_V$ 直接推出 $T$ 单射，$TR=I_U$ 直接推出 $T$ 满射。反过来，若 $T$ 单射，在 $\operatorname{Im}T$ 上定义 $L(Tx)=x$，再把 $L$ 扩张到 $U$。若 $T$ 满射，取陪域的一组基并逐一选取原像，线性扩张得到 $R$。一般地，取分解 $V=\ker T\oplus W$ 和 $U=\operatorname{Im}T\oplus Z$，则 $T|_W:W\to\operatorname{Im}T$ 是同构；定义 $G$ 在像空间上等于此同构的逆，在 $Z$ 上为零，便有 $TGT=T$。
\end{proof}

矩阵形式分别为 $BA=I_n$、$AC=I_m$ 和 $AGA=A$；其中 $A\in M_{m\times n}(K)$，$B,C,G$ 均为 $n\times m$ 矩阵。最后一种 $G$ 称为一种\emph{广义逆}（generalized inverse），通常不唯一，也不要求来自内积结构。

\begin{proposition}
\label{prop:xie-ex-factorization}
设 $T,S:V\to U$ 为线性映射，则
\[
    S=HT\text{ 对某个 }H\in\mathcal L(U)\text{ 成立}
    \quad\Longleftrightarrow\quad \ker T\subseteq\ker S,
\]
\[
    S=TF\text{ 对某个 }F\in\mathcal L(V)\text{ 成立}
    \quad\Longleftrightarrow\quad \operatorname{Im}S\subseteq\operatorname{Im}T.
\]
\end{proposition}

\begin{proof}
两个必要性分别由复合的定义得到。若 $\ker T\subseteq\ker S$，在 $\operatorname{Im}T$ 上定义 $H(Tx)=Sx$；当 $Tx=Ty$ 时，$x-y\in\ker T\subseteq\ker S$，故 $Sx=Sy$，定义良好且线性，再扩张到 $U$。若 $\operatorname{Im}S\subseteq\operatorname{Im}T$，取 $V$ 的基 $e_i$，对每个 $S(e_i)$ 选取原像 $v_i$ 使 $T(v_i)=S(e_i)$，再由 $F(e_i)=v_i$ 线性扩张，即得 $TF=S$。
\end{proof}

\subsection{幂零链与选基}

% 原书例 4.8、4.9。
\begin{lemma}
\label{lem:xie-ex-nilpotent-chain}
设 $T\in\mathcal L(V)$、$v\in V$，$q\geq1$ 为整数。若 $T^{q}(v)=0$ 而 $T^{q-1}(v)\ne0$，则 $v,Tv,\ldots,T^{q-1}v$ 线性无关，因而 $q\leq\dim V$。
\end{lemma}

\begin{proof}
若 $\sum_{j=0}^{q-1}c_jT^jv=0$ 且存在非零系数，取最小的非零系数下标 $s$，对等式作用 $T^{q-1-s}$。高于 $s$ 的项全部为零，低于 $s$ 的系数全部为零，得到 $c_sT^{q-1}v=0$，与两个因子均非零矛盾。
\end{proof}

称满足某个 $T^q=0$ 的变换为\emph{幂零变换}（nilpotent transformation）。若 $\dim V=n\geq1$，$T$ 幂零且 $\operatorname{rank}T=n-1$，则存在 $v$ 使
\[
    (v,Tv,\ldots,T^{n-1}v)
\]
成为 $V$ 的一组基，$T$ 在此基下的矩阵仅在紧邻主对角线下方有 $n-1$ 个 $1$，其余元素均为零。为证明这一点，设 $q$ 为最小幂零指数，则引理给出 $q\leq n$；另一方面，由 $\dim\ker T=1$ 以及 $T:\ker T^{j+1}\to\ker T^j$ 的核为 $\ker T$，得到 $\dim\ker T^{j+1}\leq\dim\ker T^j+1$，从而 $n=\dim\ker T^q\leq q$。因此 $q=n$，再取 $T^{n-1}v\ne0$ 即得所需基。

\section{线性同构}
\label{sec:xie-ex-linear-isomorphisms}

\subsection{插值与同构的构造}

% 原书例 4.10、4.11。
记 $K[x]_{\leq n}$ 为次数不超过 $n$ 的多项式组成的空间，包含零多项式。取互不相同的 $a_0,\ldots,a_n\in K$，定义取值映射
\[
    E:K[x]_{\leq n}\to K^{n+1},\qquad
    E(f)=(f(a_0),\ldots,f(a_n))^{\mathsf T}.
\]
这是线性映射。若 $E(f)=0$，则 $f$ 有 $n+1$ 个不同的零点；由非零的 $d$ 次多项式至多有 $d$ 个不同零点可得 $f=0$。定义域与陪域均为 $n+1$ 维，故 $E$ 为同构。于是任意给定取值 $b_0,\ldots,b_n$，都存在唯一 $f\in K[x]_{\leq n}$ 满足 $f(a_i)=b_i$。

\begin{theorem}
\label{thm:xie-ex-lagrange-interpolation}
上述唯一多项式由\emph{Lagrange 插值公式}（Lagrange interpolation formula）给出：
\begin{equation}
\label{eq:xie-ex-lagrange-interpolation}
    f(x)=\sum_{i=0}^{n}b_i\ell_i(x),\qquad
    \ell_i(x)=\prod_{\substack{0\leq j\leq n\\j\ne i}}
    \frac{x-a_j}{a_i-a_j}.
\end{equation}
\end{theorem}

\begin{proof}
由于节点互异，各分母非零；$\ell_i(a_j)$ 在 $i=j$ 时为 $1$，否则为 $0$，且 $\deg\ell_i\leq n$，故右侧多项式具有指定的全部取值。唯一性已由取值映射的单射性证明。
\end{proof}

% 原书例 4.13、4.14。
有限维等维空间之间的线性映射，只需证明单射或满射之一即可证明同构。因此一个变换可逆，当且仅当它把某一组基变成一组基；此时它也把每一组基变成一组基。进一步，若 $U,W\subseteq V$ 且 $\dim U=\dim W$，分别扩充两者的一组基到 $V$ 的基，再逐一对应基向量，就得到自同构 $A$ 使 $A(U)=W$。这里两个子空间可以相交，不要求构成直和。

\subsection{由代数关系判定可逆}

% 原书例 4.12、4.15、4.16。
\begin{proposition}
\label{prop:xie-ex-polynomial-inverse}
设 $T\in\mathcal L(V)$，$q\geq1$ 为整数，$c_0,\ldots,c_q\in K$。若
\[
    \sum_{j=0}^{q}c_jT^j=0,\qquad c_0\ne0,
\]
则
\[
    T^{-1}=-c_0^{-1}\sum_{j=1}^{q}c_jT^{j-1}.
\]
该结论也适用于无限维空间。
\end{proposition}

\begin{proof}
公式右侧是 $T$ 的多项式，与 $T$ 可交换；将已知等式的常数项移至右侧，再分别从左侧和右侧提出 $T$，便得到所给变换与 $T$ 的两个乘积均为 $I_V$。
\end{proof}

若有限维空间上的 $T$ 满足：对每个 $v$ 都存在依赖于 $v$ 的正整数 $q_v$ 使 $T^{q_v}v=0$，则取一组基并令 $q$ 为各基向量对应指数的最大值，得到统一的 $T^q=0$。于是
\begin{equation}
\label{eq:xie-ex-nilpotent-geometric-series}
    (I_V-T)^{-1}=I_V+T+\cdots+T^{q-1}.
\end{equation}
这说明局部幂零条件在有限维时推出幂零。无限维时未必有统一指数，但仍可对每个向量分别使用有限和构造 $I_V-T$ 的逆：任取使 $T^qv=0$ 的 $q$，令 $R(v)=\sum_{j=0}^{q-1}T^jv$；增大 $q$ 不改变结果，使用共同的充分大指数可验证线性以及 $(I_V-T)R=R(I_V-T)=I_V$。

\begin{example}
\label{ex:xie-ex-two-sided-multiplication}
设 $A,B\in M_n(K)$，在 $M_n(K)$ 上定义 $T(X)=AXB$。证明：$T$ 可逆当且仅当 $A,B$ 都可逆。
\end{example}

\begin{solve}
若 $A,B$ 可逆，则 $T^{-1}(Y)=A^{-1}YB^{-1}$。反之，若 $T$ 可逆，则存在 $X$ 使 $AXB=I_n$。因此
\[
    n=\operatorname{rank}(AXB)
    \leq\min\{\operatorname{rank}A,\operatorname{rank}B\}\leq n,
\]
从而 $A,B$ 均满秩且可逆。一般矩形空间上的秩与零度将在例~\ref{ex:xie-ex-rectangular-multiplication} 中计算。
\end{solve}

\subsection{无限维空间中的单侧逆与交换子}

% 原书例 4.17--4.19。
在 $K[x]$ 上定义求导 $D(f)=f'$、积分 $J(f)(x)=\int_0^x f(t)\,dt$ 和乘法变换 $M(f)=xf$。这里积分按多项式系数逐项定义，数域的特征为零保证分母可逆。直接计算得到
\[
    DJ=I,\qquad JD(f)=f-f(0),\qquad DM-MD=I.
\]
$D$ 满而不单，$J$ 单而不满，所以无限维空间中不能由一个单侧逆推出可逆。$D$ 的右逆还可取 $J_\lambda(f)=J(f)+\lambda(f)\cdot1$，其中 $\lambda:K[x]\to K$ 为任意\emph{线性泛函}（linear functional），即取值于标量域的线性映射，因而右逆一般不唯一。

\begin{proposition}
\label{prop:xie-ex-infinite-inverse-criteria}
对任意线性空间上的变换 $S,T$，两者都可逆当且仅当 $ST,TS$ 都可逆；$T$ 可逆当且仅当它恰有一个左逆，也当且仅当它恰有一个右逆。
\end{proposition}

\begin{proof}
若 $ST,TS$ 可逆，则 $(ST)^{-1}S$ 是 $T$ 的左逆，$S(TS)^{-1}$ 是 $T$ 的右逆；左逆 $L$ 与右逆 $R$ 满足 $L=L(TR)=(LT)R=R$，故 $T$ 可逆，同理 $S$ 可逆。反向由逆的乘积公式得到。再设 $LT=I$ 且左逆唯一。因 $[L+(I-TL)]T=I$，唯一性迫使 $I-TL=0$，故 $L$ 也是右逆。若 $TR=I$ 且右逆唯一，则 $T[R+(I-RT)]=I$ 同样迫使 $RT=I$。可逆变换的任何左逆或右逆都等于其逆，得到另一方向。
\end{proof}

若 $0<\dim V=n<\infty$，不可能有 $ST-TS=I_V$：在同一基下取矩阵并取迹，由 $\operatorname{tr}(AB)=\sum_{i,j}a_{ij}b_{ji}=\operatorname{tr}(BA)$，左侧迹为 $0$，右侧迹为 $n\ne0$。因此多项式空间上的 $DM-MD=I$ 也提供了有限维结论失效的具体实例。

\section{线性映射与矩阵}
\label{sec:xie-ex-linear-matrices}

\subsection{坐标同构、秩标准型与秩一分解}

% 原书例 4.20--4.25、4.28。
坐标映射 $c_{\mathcal E}:V\to K^n$ 与 $c_{\mathcal F}:U\to K^m$ 都是同构，式~\eqref{eq:xie-ex-coordinate-map} 可以写成
\[
    c_{\mathcal F}T=L_Ac_{\mathcal E},\qquad L_A(x)=Ax.
\]
因此空间问题和矩阵问题可以双向转化。相似关系也有完整的空间解释：若 $B=P^{-1}AP$，在 $K^n$ 上取 $T(x)=Ax$，标准基记为 $\mathcal E$，新基取 $\mathcal E P$，则 $T$ 在两组基下分别为 $A,B$。

\begin{proposition}
\label{prop:xie-ex-rank-normal-form}
若 $T:V\to U$ 的秩为 $r$，则可以分别选择定义域与陪域的基，使
\begin{equation}
\label{eq:xie-ex-rank-normal-form}
    [T]=\begin{pmatrix}I_r&0\\0&0\end{pmatrix}_{m\times n}.
\end{equation}
此外，$T$ 可表示为 $r$ 个秩一映射之和；$r=0$ 时理解为空和。
\end{proposition}

\begin{proof}
取 $V=\ker T\oplus W$，$W$ 的基记为 $v_1,\ldots,v_r$，则 $T(v_1),\ldots,T(v_r)$ 是像空间的基。定义域按 $W$ 的基、核的基排列，陪域按这组像、其扩充基排列，即得所述矩阵。若 $x=\sum_{i=1}^{r}a_iv_i+k$，其中 $k\in\ker T$，令 $T_i(x)=a_iT(v_i)$，则每个 $T_i$ 的像为一维，且 $T=\sum_iT_i$。
\end{proof}

\begin{notealert}[type=caution,title={秩标准型使用两侧独立换基}]
式~\eqref{eq:xie-ex-rank-normal-form} 对任意线性映射成立，即使 $V=U$，也允许定义域和陪域使用不同的基。不能据此断言每个方阵都相似于 $\operatorname{diag}(I_r,0)$；后者要求变换本身幂等，见第~\ref{sec:xie-ex-idempotents} 节。
\end{notealert}

\subsection{核相同、矩阵相抵与秩不等式}

% 原书例 4.27 与重新讨论的 Frobenius 不等式。
\begin{theorem}
\label{thm:xie-ex-same-kernel}
设 $T,S:V\to U$，则 $\ker T=\ker S$ 当且仅当存在 $U$ 上的自同构 $H$ 使 $S=HT$。因此对 $A,B\in M_{m\times n}(K)$，齐次方程组 $Ax=0$ 与 $Bx=0$ 同解，当且仅当存在 $P\in\operatorname{GL}_m(K)$ 使 $B=PA$。
\end{theorem}

\begin{proof}
若 $S=HT$ 且 $H$ 可逆，则两者的核相同。反之，在 $\operatorname{Im}T$ 上定义 $H_0(Tx)=Sx$，共同的核保证定义良好、单射且满射到 $\operatorname{Im}S$。像空间维数相同，因此可分别把这两个像空间的基扩充成 $U$ 的基，并对应扩充的基向量，把 $H_0$ 扩张为 $U$ 的自同构。矩阵结论由坐标同构得到，也说明“同解”对应初等行变换等价。
\end{proof}

若矩阵 $A\in M_{p\times q}(K)$、$B\in M_{q\times r}(K)$、$C\in M_{r\times s}(K)$，把它们视为相应空间之间的线性映射。在 $\operatorname{Im}B$ 和 $\operatorname{Im}(BC)$ 上分别限制 $A$，由维数公式有
\[
\begin{aligned}
    \operatorname{rank}(AB)
    &=\operatorname{rank}B-\dim(\ker A\cap\operatorname{Im}B),\\
    \operatorname{rank}(ABC)
    &=\operatorname{rank}(BC)-\dim(\ker A\cap\operatorname{Im}(BC)).
\end{aligned}
\]
由于 $\operatorname{Im}(BC)\subseteq\operatorname{Im}B$，相减得\emph{Frobenius 秩不等式}
\begin{equation}
\label{eq:xie-ex-frobenius-rank}
    \operatorname{rank}(ABC)\geq
    \operatorname{rank}(AB)+\operatorname{rank}(BC)-\operatorname{rank}B.
\end{equation}
证明的关键是限制映射的核，而不是只比较三个乘积的像。

\subsection{根据空间结构选择基}

% 原书例 4.26、4.29、4.30。
\begin{proposition}
\label{prop:xie-ex-basis-independent-matrix}
有限维空间上的变换 $T$ 在每一组基下的表示矩阵都相同，当且仅当 $T=cI_V$，其中 $c\in K$。
\end{proposition}

\begin{proof}
纯量变换在每组基下的矩阵均为 $cI_n$。反过来，设共同矩阵为 $A$。任取非零 $v$ 并补成基，比较该基与只把第一项变为 $2v$ 的基下的第一列，得到 $T(v)$ 在其余基方向上的分量全为零，因此每条直线都不变。固定基 $e_i$，写 $T(e_i)=c_ie_i$，再利用 $K(e_i+e_j)$ 不变可得 $c_i=c_j$，故 $T=cI_V$。零空间的情形结论亦成立。
\end{proof}

在 $M_n(K)$ 上，转置变换 $\tau(X)=X^{\mathsf T}$ 的结构可由
\[
    X=\frac{X+X^{\mathsf T}}2+\frac{X-X^{\mathsf T}}2
\]
直接看出：对称矩阵空间与反对称矩阵空间构成直和，$\tau$ 在前者上为 $I$，在后者上为 $-I$。因此取基
\[
    E_{ii},\quad E_{ij}+E_{ji}\ (i<j),\quad E_{ij}-E_{ji}\ (i<j),
\]
其中 $E_{ij}$ 为第 $(i,j)$ 元素为 $1$、其余为零的矩阵，表示矩阵便为对角阵，$1$ 与 $-1$ 分别出现 $n(n+1)/2$ 和 $n(n-1)/2$ 次。

\begin{example}
\label{ex:xie-ex-paired-square-zero}
设 $S,T\in\mathcal L(V)$ 满足 $S^2=T^2=0$、$ST+TS=I_V$。证明 $\dim V$ 为偶数，并给出两者在同一组基下的矩阵。
\end{example}

\begin{solve}
任意 $x$ 有分解 $x=STx+TSx$，其中 $STx\in\ker S$、$TSx\in\ker T$；若 $x$ 同时属于两核，则 $x=(ST+TS)x=0$。故 $V=\ker S\oplus\ker T$。对 $u\in\ker S$，有 $STu=u$，故 $T:\ker S\to\ker T$ 单射；对 $v\in\ker T$，有 $TSv=v$，故它又满射，逆为 $S|_{\ker T}$。取 $\ker S$ 的基 $u_1,\ldots,u_k$，令 $v_i=Tu_i$，则 $(u_1,v_1,\ldots,u_k,v_k)$ 是 $V$ 的基，且每一对上
\[
    [T]_i=\begin{pmatrix}0&0\\1&0\end{pmatrix},\qquad
    [S]_i=\begin{pmatrix}0&1\\0&0\end{pmatrix}.
\]
两变换的整体矩阵由这些 $2\times2$ 块组成，因而 $\dim V=2k$。
\end{solve}

\section{像空间和核空间}
\label{sec:xie-ex-images-kernels}

\subsection{从坐标还原空间中的向量}

% 原书例 4.31。
\begin{example}
\label{ex:xie-ex-compute-image-kernel}
设 $T$ 在基 $\mathcal E=(e_1,e_2,e_3,e_4)$ 下的矩阵为
\[
    A=\begin{pmatrix}
    1&0&2&1\\-1&2&1&3\\1&2&5&5\\2&-2&1&-2
    \end{pmatrix}.
\]
求 $\operatorname{Im}T$ 与 $\ker T$ 的基。
\end{example}

\begin{solve}
设列向量为 $a_1,\ldots,a_4$，则 $a_3=2a_1+\tfrac32a_2$、$a_4=a_1+2a_2$，而前两列线性无关，故秩为 $2$。将前两列还原为原空间的向量得
\[
    \operatorname{Im}T=
    \operatorname{span}\{e_1-e_2+e_3+2e_4,\ 2e_2+2e_3-2e_4\}.
\]
上述列关系给出 $Ax=0$ 的两个无关解 $(-4,-3,2,0)^{\mathsf T}$ 和 $(-1,-2,0,1)^{\mathsf T}$，由零度为 $2$ 得
\[
    \ker T=\operatorname{span}
    \{-4e_1-3e_2+2e_3,\ -e_1-2e_2+e_4\}.\qedhere
\]
\end{solve}

求像空间时，行化简可以确定主元列，但必须在\emph[caution]{原矩阵}中取对应列；行变换通常改变列空间本身。求核空间时，先求坐标解，再以原来的定义域基还原，不能把坐标列向量直接当成抽象空间中的向量。

% 原书例 4.32、4.33。
\begin{proposition}
\label{prop:xie-ex-separate-maps}
若 $T_1,\ldots,T_s:V\to U$ 都非零，则存在 $v\in V$ 使每个 $T_i(v)\ne0$。若 $S_1,\ldots,S_s:V\to U$ 两两不同，则存在 $v$ 使所有 $S_i(v)$ 两两不同。
\end{proposition}

\begin{proof}
证明第一结论只需说明无限域上的空间不能被有限个真子空间覆盖。对真子空间数目归纳：取 $a$ 不属于前 $s-1$ 个子空间，若它也不属于最后一个则已得证；否则 $a$ 属于最后一个，取 $b$ 不属于最后一个，考察直线 $a+tb$。它与前 $s-1$ 个子空间各至多交于一点，否则可由两交点相减并线性组合推出 $a$ 属于该子空间；与最后一个仅可能在 $t=0$ 处相交。数域无限，故可避开有限多个参数值。将此结论用于 $\ker T_i$ 即得第一结论，再用于全部非零差映射 $S_i-S_j$ 即得第二结论。
\end{proof}

\subsection{幂次的稳定与空间分解}

% 原书例 4.34--4.36。
对任意 $T\in\mathcal L(V)$，有子空间链
\[
    V=\operatorname{Im}T^0\supseteq\operatorname{Im}T
    \supseteq\operatorname{Im}T^2\supseteq\cdots,
    \qquad
    0=\ker T^0\subseteq\ker T\subseteq\ker T^2\subseteq\cdots.
\]

\begin{theorem}
\label{thm:xie-ex-fitting}
设 $\dim V=n$，则存在 $0\leq m\leq n$，使所有 $j\geq m$ 都满足
\[
    \operatorname{Im}T^j=\operatorname{Im}T^m,\qquad
    \ker T^j=\ker T^m,
\]
并有
\begin{equation}
\label{eq:xie-ex-fitting}
    V=\ker T^m\oplus\operatorname{Im}T^m.
\end{equation}
两部分均为 $T$ 的不变子空间，$T$ 在前者上幂零，在后者上可逆。
\end{theorem}

\begin{proof}
像的维数每次严格下降至少 $1$，故不迟于 $m=n$ 就有 $\operatorname{Im}T^m=\operatorname{Im}T^{m+1}$；对等式继续作用 $T$，所有后续像均相同。维数公式与核的包含关系使后续核也相同。若 $x=T^my\in\ker T^m$，则 $T^{2m}y=0$，由核的稳定知 $y\in\ker T^m$，故 $x=0$。两部分维数之和为 $n$，得到直和。它们的不变性由幂次可交换得到；前者上 $T^m=0$，后者上 $T$ 满射且有限维，故可逆。$m=0$ 时直和即 $V=0\oplus V$。
\end{proof}

这也称为\emph{Fitting 分解}（Fitting decomposition），并推出任意 $n$ 阶矩阵满足 $\operatorname{rank}A^n=\operatorname{rank}A^{n+1}=\cdots$。不过，通常不能把指数 $m$ 直接取成 $1$。

\begin{proposition}
\label{prop:xie-ex-first-power-splitting}
对有限维空间上的 $T$，以下条件等价：$V=\ker T\oplus\operatorname{Im}T$；$V=\ker T+\operatorname{Im}T$；$\ker T\cap\operatorname{Im}T=0$；$\ker T=\ker T^2$；$\operatorname{Im}T=\operatorname{Im}T^2$；$\ker T$ 有 $T$-不变补空间；$\operatorname{Im}T$ 有 $T$-不变补空间。核或像在一次幂与二次幂相等时，其后全部幂次也相等。
\end{proposition}

\begin{proof}
维数公式保证前两子空间的维数之和为 $n$，故“和为全空间”“交为零”“构成直和”相互等价。若 $T^2x=0$，则 $Tx\in\ker T\cap\operatorname{Im}T$；反过来，若 $Tx\in\ker T$，由两核相等可得 $x\in\ker T$。因此交为零等价于两核相等，维数公式又将其转化为两像相等，后续稳定性同定理~\ref{thm:xie-ex-fitting} 的论证。直和本身提供不变补空间。若 $V=\ker T\oplus W$ 且 $W$ 不变，则 $\operatorname{Im}T=T(W)\subseteq W$，比较维数知 $W=\operatorname{Im}T$；若 $V=\operatorname{Im}T\oplus Z$ 且 $Z$ 不变，则 $T(Z)\subseteq Z\cap\operatorname{Im}T=0$，比较维数知 $Z=\ker T$。
\end{proof}

\subsection{商空间、预先指定像核与维数估计}

% 原书例 4.37--4.44。
若 $N\subseteq V$ 为子空间，\emph{商空间}（quotient space）$V/N$ 的元素是陪集 $x+N$，运算为 $(x+N)+(y+N)=(x+y)+N$、$a(x+N)=ax+N$；若 $N$ 的基扩充为 $V$ 的基，新增基向量的陪集构成商空间的基，故 $\dim(V/N)=\dim V-\dim N$。

\begin{theorem}
\label{thm:xie-ex-first-isomorphism}
线性映射 $T:V\to U$ 诱导同构
\[
    \overline T:V/\ker T\longrightarrow\operatorname{Im}T,
    \qquad \overline T(x+\ker T)=T(x).
\]
\end{theorem}

\begin{proof}
若 $x+\ker T=y+\ker T$，则 $x-y\in\ker T$，故 $Tx=Ty$，所以定义不依赖代表元。陪集运算与 $T$ 的线性性保证 $\overline T$ 线性；每个像均有原像，而 $\overline T(x+\ker T)=0$ 当且仅当 $x\in\ker T$，即该陪集为商空间零元。因此它既满又单。
\end{proof}

取维数又得到式~\eqref{eq:xie-ex-rank-nullity}。若 $V=N\oplus W$，则 $T|_W:W\to\operatorname{Im}T$ 是同构，所以核的任意补空间均可用来选择像向量的唯一原像；特别地，满射在这样的补空间上给出到整个陪域的同构。给定 $V$ 的子空间 $N,R$，存在 $T\in\mathcal L(V)$ 使 $\ker T=N$、$\operatorname{Im}T=R$，当且仅当 $\dim N+\dim R=\dim V$：必要性来自维数公式，充分性则取补空间 $V=N\oplus W$，在 $N$ 上令 $T=0$，在 $W$ 上取任意到 $R$ 的同构。\emph[caution]{这里并不要求 $N\cap R=0$。}

作为维数公式的应用，迹映射 $\operatorname{tr}:M_n(K)\to K$ 在 $n\geq1$ 时为满射，因为 $\operatorname{tr}E_{11}=1$，故无迹矩阵空间的维数为 $n^2-1$，一组基为
\[
    \{E_{ij}:i\ne j\}\ \cup
    \{E_{ii}-E_{i+1,i+1}:1\leq i<n\}.
\]
这些矩阵的非对角元素与对角元素分别可独立比较，因而线性无关；其数目恰为所需维数。

对 $T:V\to U$、子空间 $W\subseteq V$ 与 $Z\subseteq U$，限制映射的维数公式分别给出
\begin{equation}
\label{eq:xie-ex-restriction-dimension}
    \dim T(W)=\dim W-\dim(W\cap\ker T),
\end{equation}
\begin{equation}
\label{eq:xie-ex-preimage-dimension}
    \dim T^{-1}(Z)=\dim\ker T+\dim(Z\cap\operatorname{Im}T).
\end{equation}
第二式中的 $T^{-1}(Z)$ 是集合的原像，即使 $T$ 不可逆也有意义；限制 $T:T^{-1}(Z)\to Z\cap\operatorname{Im}T$ 为满射且核仍是 $\ker T$，从而该式成立。于是
\[
    \dim W-\dim\ker T\leq\dim T(W)\leq\dim W,
    \qquad \dim T^{-1}(Z)\leq\dim\ker T+\dim Z,
\]
且当 $Z\subseteq\operatorname{Im}T$ 时第二个不等式取等号。若 $\dim V>\dim U$，维数公式还立即给出 $\ker T\ne0$。对 $A\in M_{m\times n}(K)$、$B\in M_{n\times p}(K)$，将第一式用于 $A|_{\operatorname{Im}B}$，得到\emph{Sylvester 秩不等式}
\begin{equation}
\label{eq:xie-ex-sylvester-rank}
    \operatorname{rank}A+\operatorname{rank}B-n
    \leq\operatorname{rank}(AB)
    \leq\min\{\operatorname{rank}A,\operatorname{rank}B\}.
\end{equation}

\section{不变子空间}
\label{sec:xie-ex-invariant-subspaces}

\subsection{验证不变性与可交换变换}

% 原书例 4.45--4.50。
若 $W=\operatorname{span}\{w_1,\ldots,w_s\}$，验证 $W$ 不变只需逐一验证 $T(w_i)\in W$，不要求生成向量线性无关。两个不变子空间的交与和仍不变，因为像既保持各自归属，又保持向量的加法分解；相同论证也适用于有限个交与和。每条一维子空间都不变，当且仅当 $T$ 为纯量变换，证明见命题~\ref{prop:xie-ex-basis-independent-matrix} 中对 $e_i$ 与 $e_i+e_j$ 的比较。

更一般地，对固定的 $1\leq r\leq n-1$，如果每个 $r$ 维子空间都不变，那么每条直线也不变。事实上，给定直线 $L$，所有包含 $L$ 的 $r$ 维子空间之交就是 $L$：对任何 $x\notin L$，扩充 $L$ 的基并把 $x$ 作为第二个基向量，可以选出包含 $L$ 而避开 $x$ 的 $r$ 维坐标子空间。因此 $T$ 保持上述交，从而是纯量变换。反过来纯量变换保持所有子空间。这就得到“任意固定非平凡维数的全部子空间不变”与“纯量变换”等价。

\begin{proposition}
\label{prop:xie-ex-commuting-invariance}
若 $ST=TS$，则 $\ker S$ 和 $\operatorname{Im}S$ 都是 $T$ 的不变子空间。
\end{proposition}

\begin{proof}
若 $Sx=0$，则 $S(Tx)=T(Sx)=0$；若 $x=Sy$，则 $Tx=TSy=STy\in\operatorname{Im}S$。
\end{proof}

例如，设 $A$ 为幂零矩阵、$AB=BA$，且 $\operatorname{rank}(AB)=\operatorname{rank}B$。像空间 $\operatorname{Im}B$ 对 $A$ 不变，且 $A$ 在其上的限制满射；若 $A^q=0$，这个满射的 $q$ 次复合仍满射，却同时是零映射，只能有 $\operatorname{Im}B=0$，故 $B=0$。

若 $T$ 可逆且 $W$ 是有限维 $T$-不变子空间，则 $T|_W$ 单射，因而满射，故 $T(W)=W$、$T^{-1}(W)=W$。有限维条件不能随意去掉：在 Laurent 多项式空间 $K[x,x^{-1}]$ 上，乘以 $x$ 为自同构，$K[x]$ 对它不变，但乘以 $x^{-1}$ 并不保持 $K[x]$。

\subsection{求导变换与互异对角元的全部不变子空间}

% 原书例 4.51、4.52。
\begin{example}
\label{ex:xie-ex-derivative-invariant}
设 $n\geq1$。在 $K[x]_{<n}=\{f:\deg f<n\}\cup\{0\}$ 上，令 $D(f)=f'$。求 $D$ 的全部不变子空间。
\end{example}

\begin{solve}
除零空间外，取不变子空间 $W$ 中次数最大的多项式 $f$，记次数为 $d$。$f,Df,\ldots,D^df$ 均在 $W$ 中，次数分别为 $d,d-1,\ldots,0$，最高项系数均非零，故线性无关并张成 $K[x]_{\leq d}$；按 $d$ 的最大性又有 $W\subseteq K[x]_{\leq d}$。因此 $W=K[x]_{\leq d}$。全部不变子空间恰为
\[
    0,\quad K[x]_{\leq0},\quad K[x]_{\leq1},\quad\ldots,
    \quad K[x]_{\leq n-1}.
\]
特别地，每个非零不变子空间都包含常数多项式 $1$。
\end{solve}

\begin{example}
\label{ex:xie-ex-distinct-diagonal-invariant}
设 $T(e_i)=d_ie_i$，其中 $d_1,\ldots,d_n\in K$ 两两不同。求 $T$ 的全部不变子空间。
\end{example}

\begin{solve}
定义多项式
\[
    p_i(t)=\prod_{j\ne i}\frac{t-d_j}{d_i-d_j}.
\]
于是 $p_i(T)e_j=0$（$j\ne i$），$p_i(T)e_i=e_i$，即 $p_i(T)$ 提取第 $i$ 个坐标分量。若 $W$ 对 $T$ 不变，则它对 $T$ 的所有多项式也不变；只要 $w=\sum_jc_je_j\in W$ 且 $c_i\ne0$，就有 $c_ie_i=p_i(T)w\in W$，从而 $e_i\in W$。故 $W$ 必由某个基向量子集张成，反过来每个这样的子空间都不变，共有 $2^n$ 个，包含零空间与全空间。互异条件保证分母非零；若所有 $d_i$ 相同，则每个子空间均不变，结论会发生根本变化。
\end{solve}

\subsection{限制变换与商空间上的诱导变换}

% 原书例 4.53。
若 $W$ 为 $T$-不变子空间，则不仅能研究 $T|_W$，还可在商空间上定义
\[
    \widetilde T:V/W\to V/W,\qquad
    \widetilde T(x+W)=T(x)+W.
\]
当 $x-y\in W$ 时，$T(x)-T(y)\in W$，所以定义良好；陪集运算直接保证其线性。对式~\eqref{eq:xie-ex-invariant-block} 中的适配基，$T|_W$ 的矩阵为 $A_{11}$，$\widetilde T$ 在商空间基 $e_{r+1}+W,\ldots,e_n+W$ 下的矩阵为 $A_{22}$。非对角块 $A_{12}$ 描述所选补空间中的向量映入 $W$ 的分量；只有补空间也不变时，它才为零，不能因已有一个不变子空间就直接写成分块对角阵。

\section{幂等变换}
\label{sec:xie-ex-idempotents}

\subsection{幂等变换就是投影}

% 原书例 4.54--4.56。
\begin{definition}
满足 $P^2=P$ 的线性变换称为\emph{幂等变换}（idempotent transformation）。若 $V=U\oplus W$，将 $u+w$ 映为 $u$ 的变换称为沿 $W$ 到 $U$ 的\emph{投影}（projection）。这里不涉及内积，因而不要求正交。
\end{definition}

\begin{theorem}
\label{thm:xie-ex-idempotent-projection}
$P^2=P$ 当且仅当 $P$ 是沿其核到其像的投影。此时
\[
    V=\operatorname{Im}P\oplus\ker P,\qquad
    \operatorname{Im}P=\ker(I-P),\qquad
    \ker P=\operatorname{Im}(I-P).
\]
若 $\operatorname{rank}P=r$，则存在一组基使 $[P]=\operatorname{diag}(I_r,0)$。
\end{theorem}

\begin{proof}
若 $P^2=P$，则 $x=Px+(I-P)x$，其中 $P(Px)=Px$、$P((I-P)x)=0$。若 $y=Px\in\ker P$，则 $y=P^2x=Py=0$，因此构成直和。$P$ 在其像上为恒等，在核上为零，正是所述投影；两个像核等式可由上述分解逐项验证。反过来，投影作用一次后已经在目标子空间中，再作用不变，故 $P^2=P$。拼接像与核的基即得矩阵形式。
\end{proof}

因此幂等矩阵的秩等于其迹，两个同阶幂等矩阵相似当且仅当秩相同。特别地，若 $A,B$ 幂等且同秩，则存在可逆矩阵 $C$ 使 $B=C^{-1}AC$，即 $CB=AC$。这些结论来自投影的空间分解，无须预先使用特征值理论。

\subsection{两个投影的像、核与加减}

% 原书例 4.57、4.58。
\begin{proposition}
\label{prop:xie-ex-two-projections}
设 $P^2=P$、$Q^2=Q$，则
\[
\begin{aligned}
    \operatorname{Im}P=\operatorname{Im}Q
    &\iff PQ=Q,\quad QP=P,\\
    \ker P=\ker Q
    &\iff PQ=P,\quad QP=Q,\\
    (P+Q)^2=P+Q
    &\iff PQ=QP=0,\\
    (P-Q)^2=P-Q
    &\iff PQ=QP=Q.
\end{aligned}
\]
\end{proposition}

\begin{proof}
投影在自身像上为恒等，因此两像相同推出第一组乘积等式，反过来乘积等式给出两个方向的像包含。若两核相同，利用 $x-Qx\in\ker Q=\ker P$ 得 $Px=PQx$，同理 $Q=QP$；反向则由 $Px=0\Rightarrow Qx=QPx=0$ 及对称论证得到。

对和的幂等性，展开得 $PQ+QP=0$。分别从左、右乘 $P$，得到 $PQ+PQP=0$、$PQP+QP=0$，所以 $PQ=QP$，再由 $2PQ=0$ 得二者为零。对差，展开得 $PQ+QP=2Q$。左右乘 $P$ 分别得到 $PQP=PQ$ 和 $PQP=QP$，故 $PQ=QP$，再代回得 $PQ=Q$。各充分性均由展开验证；这里使用了数域中特征不为 $2$ 的事实。
\end{proof}

\subsection{多个投影、直和与秩的可加性}

% 原书例 4.59、4.60，以及训练题三 10、11。
\begin{theorem}
\label{thm:xie-ex-family-projections}
设 $P_1,\ldots,P_s\in\mathcal L(V)$ 满足 $P_i^2=P_i$、$P_iP_j=0$（$i\ne j$）。令 $P=\sum_iP_i$，则 $P^2=P$，并且
\begin{equation}
\label{eq:xie-ex-family-projections}
    V=\operatorname{Im}P_1\oplus\cdots\oplus\operatorname{Im}P_s
    \oplus\bigcap_{i=1}^{s}\ker P_i.
\end{equation}
因此 $\bigcap_i\ker P_i=0$ 当且仅当 $P=I_V$，此时 $P_i$ 正是该直和分解的各分量投影。
\end{theorem}

\begin{proof}
展开乘积得 $P^2=P$，且 $P_iP=PP_i=P_i$。若 $x_i\in\operatorname{Im}P_i$ 且 $\sum_ix_i=0$，作用 $P_j$ 得 $x_j=0$，故各像的和为直和，其和等于 $\operatorname{Im}P$。又 $Px=0$ 推出 $P_ix=P_iPx=0$，反向由求和得到，所以 $\ker P=\bigcap_i\ker P_i$。应用定理~\ref{thm:xie-ex-idempotent-projection} 即得分解及其余结论。
\end{proof}

\begin{theorem}
\label{thm:xie-ex-idempotent-rank-additivity}
设 $P^2=P$ 且 $P=T_1+\cdots+T_s$，则
\begin{equation}
\label{eq:xie-ex-idempotent-rank-additivity}
    \operatorname{rank}P=\sum_{i=1}^{s}\operatorname{rank}T_i
    \quad\Longleftrightarrow\quad
    T_i^2=T_i,\quad T_iT_j=0\ (i\ne j).
\end{equation}
\end{theorem}

\begin{proof}
若右侧成立，定理~\ref{thm:xie-ex-family-projections} 给出像空间的直和，因而秩可加。反过来，设左侧成立，记 $R_i=\operatorname{Im}T_i$。由
\[
    \operatorname{Im}P\subseteq\sum_iR_i,\qquad
    \dim\operatorname{Im}P=\sum_i\dim R_i
\]
以及和空间维数不超过各维数之和，得 $\operatorname{Im}P=\bigoplus_iR_i$。对任意 $y\in R_i$，因 $P$ 在其像上为恒等，$y=Py=\sum_jT_jy$；由直和表示的唯一性，$T_iy=y$、$T_jy=0$（$j\ne i$）。令 $y=T_ix$ 即得 $T_i^2=T_i$、$T_jT_i=0$。
\end{proof}

特别地，若 $\sum_iT_i=I_V$ 且 $\sum_i\operatorname{rank}T_i=\dim V$，则各 $T_i$ 自动成为互相消去的投影。几何上，秩的可加性排除了各像空间之间的重叠，而和为恒等又保证这些像覆盖整个空间。

\section{基础训练：方法与综合应用}
\label{sec:xie-ex-linear-training}

\subsection{基本判断与换基计算}

基础判断先归结为定义与维数，而不直接套选项。含非零常数平移的映射不保持零向量；坐标中出现平方的映射一般不保持数乘，例如 $(x_1,x_2,x_3)\mapsto(x_1^2,x_2,x_3)$ 在 $x_1=1$ 处与数乘 $2$ 不相容。一维空间上的每个线性变换都是 $cI$，$n$ 维空间的变换空间维数为 $n^2$。两个变换在一组基上的作用相同便相等，但仅有相同的秩、像或核均不能推出相等；例如 $I$ 与 $2I$ 有相同的像和核。线性相关组的像仍相关，因为原有非零系数关系在作用线性映射后仍成立。

维数公式并不自动给出像与核的直和。例如在 $K[x]_{<2}$ 上，求导 $D$ 的像与核均为 $K\cdot1$，其维数之和虽为 $2$，交却非零。幂等变换保证像核直和；幂零变换则可能有非零的像核交。非零有限维变换若没有非平凡不变子空间，则其核只能为零，从而可逆。对实二维空间，旋转 $\pi/2$ 的变换 $R(x,y)=(-y,x)$ 没有不变直线：若 $R(v)=cv$ 且 $v\ne0$，由 $R^2=-I$ 得 $c^2=-1$，与 $c\in\mathbb R$ 矛盾。这也说明不变子空间问题必须明确数域。

计算时坚持使用“基向量的像为列”的约定。例如 $D$ 在 $1,x,\ldots,x^{n-1}$ 下的矩阵满足第 $(i,i+1)$ 元素为 $i$（$1\leq i<n$），其余为零。若 $S,T$ 在同一组基下的矩阵为 $A,B$，则 $ST+2S^2$ 的矩阵为 $AB+2A^2$；若 $S$ 可逆，基的过渡矩阵为 $P$，则 $TS^{-2}+2S+I$ 在新基下的矩阵为
\[
    P^{-1}(BA^{-2}+2A+I_n)P.
\]
若恰有 $T(e_i)=f_i$ 且 $(f_1,\ldots,f_n)=(e_1,\ldots,e_n)P$，则 $T$ 在旧基下为 $P$，在新基下仍为 $P^{-1}PP=P$；这是该映射的特殊关系，不能推广为一般换基不改变矩阵。

\subsection{构造、幂零链与复合映射}

% 对应训练题三 1--7；已在正文完整证明的结论不重复展开。
\begin{example}
\label{ex:xie-ex-map-space-basis}
取 $V$ 的基 $e_1,\ldots,e_n$ 和 $U$ 的基 $f_1,\ldots,f_m$，定义 $T_{ij}(e_j)=f_i$、$T_{ij}(e_k)=0$（$k\ne j$）。证明这些映射构成 $\mathcal L(V,U)$ 的一组基。
\end{example}

\begin{solve}
若 $T(e_j)=\sum_i a_{ij}f_i$，则在线性扩张定理下 $T=\sum_{i,j}a_{ij}T_{ij}$，故这些映射张成整个空间。若 $\sum_{i,j}c_{ij}T_{ij}=0$，将等式作用于 $e_j$ 得 $\sum_ic_{ij}f_i=0$，由 $f_i$ 线性无关可知全部 $c_{ij}=0$。因此共有 $mn$ 个基向量，且对应 $T_{ij}\mapsto E_{ij}$ 正是表示矩阵同构。
\end{solve}

有限支撑数列空间 $K^{(\mathbb N_0)}$ 是所有仅有有限多个非零项的数列组成的空间，映射 $(a_0,a_1,\ldots)\mapsto\sum_{j\geq0}a_jx^j$ 是到 $K[x]$ 的线性同构：该和实际有限，系数的唯一性保证单射，每个多项式又都来自其系数数列。对 $K[x]$ 上的 $D(f)=f'$、$M(f)=xf$，乘积求导对任意正整数 $q$ 进一步给出
\[
    DM^q-M^qD=qM^{q-1}\qquad(q\geq1),
\]
因为 $D(x^qf)-x^qD(f)=qx^{q-1}f$。这推广了前面的交换子等式。

若 $V=U_1\oplus U_2$，且 $T(U_1)\subseteq U_2$、$T(U_2)\subseteq U_1$，则 $\operatorname{Im}T=T(U_1)\oplus T(U_2)$，因为两像分别落在交为零的两个子空间中。因此 $\operatorname{rank}T=\dim T(U_1)+\dim T(U_2)$，该结论不要求 $U_1,U_2$ 维数相等。

\begin{example}
\label{ex:xie-ex-cyclic-nilpotent-invariant}
设 $T(e_i)=e_{i+1}$（$i<n$）、$T(e_n)=0$。证明：包含 $e_1$ 的不变子空间只能是 $V$；每个非零不变子空间都包含 $e_n$；$V$ 不能分解为两个非零不变子空间的直和。
\end{example}

\begin{solve}
包含 $e_1$ 且不变就必须包含 $Te_1=e_2$ 及其后所有基向量，因此等于 $V$。若 $W\ne0$ 不变，取 $w=\sum_i a_ie_i\ne0$，令 $k$ 为最小非零系数下标，则 $T^{n-k}w=a_ke_n\in W$，故 $e_n\in W$。若有两个非零不变子空间构成直和，它们都包含 $e_n$，交非零，产生矛盾。
\end{solve}

若 $A$ 可逆，则 $BA=A^{-1}(AB)A$，所以 $AB$ 与 $BA$ 相似；这里 $B$ 无须可逆。对一般复合映射 $ST$，把 $T$ 限制到 $\ker(ST)$，其像包含于 $\ker S$，核正是 $\ker T$，由维数公式有
\begin{equation}
\label{eq:xie-ex-product-nullity}
    \dim\ker(ST)\leq\dim\ker S+\dim\ker T.
\end{equation}
对同一 $n$ 维空间上的变换，这也等价于式~\eqref{eq:xie-ex-sylvester-rank} 的下界。

\subsection{限制映射、对合与矩形矩阵空间}

% 对应训练题三 8--13；三 10、11 已归入投影族定理。
\begin{example}
\label{ex:xie-ex-restricted-matrix-image}
设 $A\in M_{m\times n}(K)$、$B\in M_{\ell\times n}(K)$，$W=\ker B$，$T(x)=Ax$。证明
\[
    \dim T(W)=\operatorname{rank}\begin{pmatrix}A\\B\end{pmatrix}
    -\operatorname{rank}B.
\]
\end{example}

\begin{solve}
限制 $T|_W$ 的核为 $\ker A\cap\ker B$，等于纵向拼接矩阵的核。又 $\dim W=n-\operatorname{rank}B$，故由式~\eqref{eq:xie-ex-restriction-dimension} 得
\[
    \dim T(W)=(n-\operatorname{rank}B)
    -\left(n-\operatorname{rank}\begin{pmatrix}A\\B\end{pmatrix}\right),
\]
化简即得结论。
\end{solve}

\begin{example}
\label{ex:xie-ex-involution}
设 $T^2=I_V$，称 $T$ 为\emph{对合变换}（involution）。证明
\[
    V=\ker(T-I_V)\oplus\ker(T+I_V).
\]
\end{example}

\begin{solve}
任意 $v$ 可写为
\[
    v=\frac{v+Tv}{2}+\frac{v-Tv}{2}.
\]
由 $T^2=I_V$，第一项被 $T$ 保持，第二项被 $T$ 映为其相反向量。若 $Tv=v=-Tv$，则 $2v=0$，故交为零。相应分量投影是 $(I_V+T)/2$ 和 $(I_V-T)/2$。条件 $T\ne I_V$ 只保证第二个子空间非零，不保证第一个也非零，例如 $T=-I_V$。
\end{solve}

\begin{example}
\label{ex:xie-ex-rational-cyclic}
设 $V$ 是有理数域上的三维空间，$T(\alpha)=\beta$、$T(\beta)=\gamma$、$T(\gamma)=\alpha+\beta$，且 $\alpha\ne0$。证明 $\alpha,\beta,\gamma$ 线性无关。
\end{example}

\begin{solve}
先说明两个将用到的算术事实：多项式 $t^3-t-1$ 和 $t^3-t+1$ 都没有有理根。事实上，若约分后的 $a/b$ 是其中某个多项式的根，由整系数首一方程可知 $b=1$ 且 $a$ 整除常数项，所以只需检验 $a=1,-1$，代入均不为零。

若 $\alpha,\beta$ 相关，由 $\alpha\ne0$ 可设 $\beta=c\alpha$，其中 $c\in\mathbb Q$。于是 $\gamma=c^2\alpha$，再由 $T\gamma=\alpha+\beta$ 得 $c^3=c+1$，与第一个算术事实矛盾。因此 $\alpha,\beta$ 无关。若三向量相关，便可设 $\gamma=u\alpha+v\beta$，其中 $u,v\in\mathbb Q$；作用 $T$ 并与 $T\gamma=\alpha+\beta$ 比较，得到
\[
    uv=1,\qquad u+v^2=1.
\]
因此 $v\ne0$，消去 $u$ 得 $v^3-v+1=0$，与第二个算术事实矛盾。故三向量线性无关；有理数域条件在排除这些方程的根时使用，不能省略。
\end{solve}

\begin{example}
\label{ex:xie-ex-rectangular-multiplication}
设 $A\in M_m(K)$、$B\in M_n(K)$，定义 $T:M_{m\times n}(K)\to M_{m\times n}(K)$ 为 $T(X)=AXB$。求 $T$ 的秩、零度及核的一组基。
\end{example}

\begin{solve}
记 $r=\operatorname{rank}A$、$s=\operatorname{rank}B$，取可逆矩阵 $L_A,R_A,L_B,R_B$ 使
\[
    L_AAR_A=\operatorname{diag}(I_r,0),\qquad
    L_BBR_B=\operatorname{diag}(I_s,0).
\]
对 $X=R_AYL_B$，有
\[
    L_AT(X)R_B=
    \operatorname{diag}(I_r,0)\,Y\,\operatorname{diag}(I_s,0).
\]
左右两侧的坐标变换都可逆，而右侧仅保留 $Y$ 的前 $r$ 行、前 $s$ 列组成的块。因此
\[
    \operatorname{rank}T=rs,\qquad \dim\ker T=mn-rs,
\]
且核的一组基为
\[
    \{R_AE_{ij}L_B:1\leq i\leq m,\ 1\leq j\leq n,
    \ i>r\text{ 或 }j>s\}.
\]
这里 $E_{ij}$ 是 $m\times n$ 基础矩阵；上述条件恰好去除左上角 $r\times s$ 块。基向量的独立性与完备性由可逆变换 $Y\mapsto R_AYL_B$ 保证。
\end{solve}

\subsection{线性关系的保持与秩可加的同时标准型}

% 对应训练题三 14、15。
\begin{example}
\label{ex:xie-ex-relations-preservation}
给定 $n$ 维空间 $V$ 中的两个向量组 $S=(v_1,\ldots,v_m)$、$R=(w_1,\ldots,w_m)$，定义关系空间
\[
    \mathcal R_S=\{(a_1,\ldots,a_m)^{\mathsf T}\in K^m:
    \textstyle\sum_i a_iv_i=0\},
\]
并类似定义 $\mathcal R_R$。证明：存在 $T\in\mathcal L(V)$ 使 $T(v_i)=w_i$ 当且仅当 $\mathcal R_S\subseteq\mathcal R_R$；存在满足此条件的自同构，当且仅当 $\mathcal R_S=\mathcal R_R$。
\end{example}

\begin{solve}
关系空间是线性映射 $(a_i)\mapsto\sum_i a_iv_i$ 的核，所以是子空间。若存在 $T$，对任何线性关系作用 $T$，立即得到关系空间包含；若 $T$ 可逆，对逆映射同样论证，得到相等。反过来，若包含成立，在 $\operatorname{span}S$ 上定义
\[
    T_0\left(\sum_i a_iv_i\right)=\sum_i a_iw_i.
\]
两种表示之差属于 $\mathcal R_S$，从而也属于 $\mathcal R_R$，所以定义良好且线性。扩张 $T_0$ 到 $V$ 即得所需变换。若两关系空间相等，则此定义还给出 $\operatorname{span}S\to\operatorname{span}R$ 的同构；分别扩充两边的基为 $V$ 的基，再对应补充的基向量，即扩张为自同构。
\end{solve}

\begin{example}
\label{ex:xie-ex-additive-rank-normal-form}
设 $A,B\in M_{m\times n}(K)$，且 $\operatorname{rank}(A+B)=\operatorname{rank}A+\operatorname{rank}B$。记 $r=\operatorname{rank}A$、$s=\operatorname{rank}B$。证明存在可逆矩阵 $P\in\operatorname{GL}_m(K)$、$Q\in\operatorname{GL}_n(K)$，使
\[
    PAQ=\begin{pmatrix}I_r&0&0\\0&0&0\\0&0&0\end{pmatrix},
    \qquad
    PBQ=\begin{pmatrix}0&0&0\\0&I_s&0\\0&0&0\end{pmatrix},
\]
两矩阵的行分块大小均为 $r,s,m-r-s$，列分块大小均为 $r,s,n-r-s$。
\end{example}

\begin{solve}
将 $A,B$ 看作 $K^n\to K^m$ 的映射。由
\[
    \operatorname{Im}(A+B)\subseteq\operatorname{Im}A+\operatorname{Im}B,
    \qquad \operatorname{rank}(A+B)=r+s
\]
及和空间维数上界，得到 $\operatorname{Im}(A+B)=\operatorname{Im}A\oplus\operatorname{Im}B$。另一方面，纵向拼接矩阵满足
\[
    r+s=\operatorname{rank}(A+B)
    \leq\operatorname{rank}\begin{pmatrix}A\\B\end{pmatrix}
    \leq r+s,
\]
故 $N=\ker A\cap\ker B$ 的维数为 $n-r-s$，且由交和维数公式有 $\ker A+\ker B=K^n$。

取 $N$ 在 $\ker B$ 中的补空间 $X$、在 $\ker A$ 中的补空间 $Y$，则 $\dim X=r$、$\dim Y=s$，并有 $K^n=X\oplus Y\oplus N$。取基 $x_1,\ldots,x_r$、$y_1,\ldots,y_s$、$z_1,\ldots,z_{n-r-s}$，则 $Ax_i$ 是 $\operatorname{Im}A$ 的基，$By_j$ 是 $\operatorname{Im}B$ 的基；两组像共同线性无关，故可扩充为 $K^m$ 的基。在这些定义域与陪域基下，两映射恰有题述矩阵形式。若 $Q$ 的列为定义域新基、$F$ 的列为陪域新基，则取 $P=F^{-1}$ 即得结论。
\end{solve}

% 教材第 4 章（高代习题-4.pdf）整理内容至此结束。

% BEGIN textbook-xie-exercises-ch5




